\documentclass[12pt, a4paper]{article}

% ── Packages ──────────────────────────────────────────────────────
\usepackage{fontspec}
\setmainfont{Arial}
\usepackage{amsmath, amssymb, amsthm}
\usepackage{mathtools}
\usepackage{enumitem}
\usepackage{tikz}
\usetikzlibrary{calc, positioning, arrows.meta, fit, patterns, decorations.pathreplacing}
\usepackage{geometry}
\geometry{margin=2.4cm}
\usepackage{xcolor}
\usepackage{tcolorbox}
\tcbuselibrary{theorems, skins, breakable}
\usepackage{booktabs}
\usepackage{hyperref}
\hypersetup{colorlinks=true, linkcolor=blue!60!black, urlcolor=blue!60!black}

\setlength{\parindent}{0pt}
\setlength{\parskip}{6pt}

% ── Colours ───────────────────────────────────────────────────────
\definecolor{setA}{RGB}{66, 133, 244}   % blue
\definecolor{setB}{RGB}{234, 67, 53}    % red
\definecolor{setC}{RGB}{251, 188, 4}    % yellow
\definecolor{theoryBg}{RGB}{232, 245, 253}
\definecolor{theoryFrame}{RGB}{66, 133, 244}
\definecolor{stmtBg}{RGB}{255, 243, 224}
\definecolor{stmtFrame}{RGB}{230, 126, 34}
\definecolor{ansBg}{RGB}{232, 250, 232}
\definecolor{ansFrame}{RGB}{39, 174, 96}

% ── Environments ──────────────────────────────────────────────────
\newtcolorbox{theorybox}[1][]{%
  colback=theoryBg, colframe=theoryFrame,
  fonttitle=\bfseries, title={\small Տեսության ամփոփում. #1},
  breakable, enhanced, arc=2mm,
  left=5pt, right=5pt, top=5pt, bottom=5pt
}

\newtcolorbox{problembox}{%
  colback=stmtBg, colframe=stmtFrame,
  fonttitle=\bfseries, title={\small Խնդրի դրվածք},
  breakable, enhanced, arc=2mm,
  left=5pt, right=5pt, top=5pt, bottom=5pt
}

\newtcolorbox{answerbox}{%
  colback=ansBg, colframe=ansFrame,
  arc=2mm, left=5pt, right=5pt, top=3pt, bottom=3pt
}

\newcommand{\Z}{\mathbb{Z}}
\newcommand{\R}{\mathbb{R}}
\newcommand{\N}{\mathbb{N}}
\newcommand{\powset}[1]{2^{#1}}
\newcommand{\comp}[1]{\overline{#1}}
\newcommand{\symdiff}{\mathbin{\triangle}}

% ── Title ─────────────────────────────────────────────────────────
\title{\textbf{Գործնական 8 — Անդրադարձ առնչություններ}\\[4pt]
       \large Ընտրված խնդիրների լուծումներ\\[2pt]
       \normalsize Խնդիրներ 2, 6}
\author{Դիսկրետ Մաթեմատիկա — ԵՊՀ}
\date{}

\begin{document}
\maketitle
\tableofcontents
\newpage

% ══════════════════════════════════════════════════════════════════
\section{Խնդիր 2 — Երկրորդ կարգի գծային անդրադարձ առնչության լուծում}
% ══════════════════════════════════════════════════════════════════

\begin{theorybox}[Հաստատուն գործակիցներով երկրորդ կարգի գծային անդրադարձ առնչություններ]
\textbf{Հաստատուն գործակիցներով երկրորդ կարգի գծային անդրադարձ առնչությունն} ունի հետևյալ տեսքը՝
\[
  a_n = c_1\, a_{n-1} + c_2\, a_{n-2}, \qquad n \ge 2,
\]
որտեղ $c_1, c_2$ հաստատուններ են, իսկ $a_0, a_1$ (կամ $a_1, a_2$) տրված սկզբնական արժեքներն են։

\medskip
\textbf{Բնութագրիչ հավասարման մեթոդ։}\;
Առնչության մեջ տեղադրելով $a_n = r^n$ փորձնական լուծումը՝ կստանանք \textbf{բնութագրիչ հավասարումը}։
\[
  r^2 - c_1\, r - c_2 = 0.
\]

Ընդհանուր լուծման տեսքը կախված է $r_1, r_2$ արմատներից։

\medskip
\begin{center}
\renewcommand{\arraystretch}{1.5}
\begin{tabular}{lll}
\toprule
\textbf{Դեպք} & \textbf{Արմատներ} & \textbf{Ընդհանուր Լուծում} \\
\midrule
Տարբեր իրական արմատներ & $r_1 \neq r_2$, երկուսն էլ իրական
  & $a_n = \alpha\, r_1^n + \beta\, r_2^n$ \\
Պատիկ իրական արմատներ & $r_1 = r_2 = r$
  & $a_n = (\alpha + \beta\, n)\, r^n$ \\
Կոմպլեքս համալուծ արմատներ & $r_{1,2} = \rho\, e^{\pm i\theta}$
  & $a_n = \rho^n\bigl(\alpha \cos n\theta + \beta \sin n\theta\bigr)$ \\
\bottomrule
\end{tabular}
\end{center}

\medskip
$\alpha$ և $\beta$ հաստատունները որոշվում են ընդհանուր լուծման մեջ սկզբնական պայմանները տեղադրելու և ստացված $2 \times 2$ գծային հավասարումների համակարգը լուծելու միջոցով։

\begin{center}
\begin{tikzpicture}[>=Stealth, font=\small, node distance=1.6cm]
  \node[draw, rounded corners=4pt, fill=stmtBg, minimum width=2.8cm,
        align=center, font=\small] (rec) {Առնչություն\\$a_n = c_1 a_{n-1} + c_2 a_{n-2}$};
  \node[draw, rounded corners=4pt, fill=theoryBg, minimum width=2.8cm,
        align=center, font=\small, right=2.6cm of rec] (char)
    {Բնութագրիչ հավ.\\$r^2 - c_1 r - c_2 = 0$};
  \node[draw, rounded corners=4pt, fill=setC!15, minimum width=2.8cm,
        align=center, font=\small, right=2.6cm of char] (gen)
    {Ընդհանուր լուծում\\$a_n = \alpha r_1^n + \beta r_2^n$};
  \node[draw, rounded corners=4pt, fill=ansBg, minimum width=2.8cm,
        align=center, font=\small, below=1.4cm of gen] (sol)
    {Մասնավոր լուծում\\(սկզբն. պայմաններ)};

  \draw[->, thick] (rec) -- (char)
    node[midway, above=3pt, font=\scriptsize]{տեղադրել $a_n = r^n$};
  \draw[->, thick] (char) -- (gen)
    node[midway, above=3pt, font=\scriptsize]{գտնել արմատները};
  \draw[->, thick] (gen) -- (sol)
    node[midway, right=3pt, font=\scriptsize]{գտնել $\alpha, \beta$};
\end{tikzpicture}
\end{center}
\end{theorybox}

\begin{problembox}
Լուծել գծային անդրադարձ առնչությունը.
\[
  a_n = 5a_{n-1} - 6a_{n-2}, \qquad n \ge 3,
\]
տրված $a_1 = 2$ և $a_2 = 10$ սկզբնական պայմաններով։
\end{problembox}

\subsection*{Լուծում}

\textbf{Ո՞ր դեպքն է։}\; Բնութագրիչ հավասարման արմատները գտնելուց հետո ընտրում ենք ընդհանուր լուծման ձևը։
\begin{itemize}[nosep]
  \item Տարբեր իրական արմատներ $r_1 \neq r_2$. $a_n = \alpha\, r_1^n + \beta\, r_2^n$։
  \item Պատիկ իրական արմատներ $r_1 = r_2 = r$. $a_n = (\alpha + \beta\, n)\, r^n$։
  \item Կոմպլեքս համալուծ արմատներ $\rho\, e^{\pm i\theta}$. $a_n = \rho^n(\alpha \cos n\theta + \beta \sin n\theta)$։
\end{itemize}

\medskip
\textbf{Քայլ 1. Գրել բնութագրիչ հավասարումը։}

$a_n - 5a_{n-1} + 6a_{n-2} = 0$ առնչության բնութագրիչ հավասարումն է՝
\[
  r^2 - 5r + 6 = 0.
\]

\textbf{Քայլ 2. Գտնել արմատները։}

Վերլուծելով արտադրիչների՝
\[
  r^2 - 5r + 6 = (r - 2)(r - 3) = 0
  \qquad\Longrightarrow\qquad r_1 = 2, \quad r_2 = 3.
\]
Քանի որ արմատներն իրական են և իրարից տարբեր, մենք գործ ունենք \textbf{առաջին դեպքի} հետ։

\begin{center}
\begin{tikzpicture}[>=Stealth, font=\small]
  % Parabola for r^2 - 5r + 6
  \draw[->, thick] (-0.5, 0) -- (5.5, 0) node[right]{$r$};
  \draw[->, thick] (0, -1.5) -- (0, 4) node[above]{$r^2 - 5r + 6$};

  \draw[thick, setA, domain=0.5:4.5, samples=60]
    plot (\x, {(\x - 2)*(\x - 3)});

  % Roots
  \fill[setB] (2, 0) circle (3pt) node[below left=3pt, font=\footnotesize]{$r_1 = 2$};
  \fill[setB] (3, 0) circle (3pt) node[below right=3pt, font=\footnotesize]{$r_2 = 3$};

  % Tick marks
  \foreach \x in {1, 2, 3, 4} {
    \draw (\x, 0.08) -- (\x, -0.08);
  }
  \node[below, font=\tiny] at (1, -0.08) {$1$};
  \node[below, font=\tiny] at (4, -0.08) {$4$};
\end{tikzpicture}
\end{center}

\textbf{Քայլ 3. Գրել ընդհանուր լուծումը։}

Քանի որ $r_1 = 2$ և $r_2 = 3$ տարբեր իրական արմատներ են՝
\[
  a_n = \alpha \cdot 2^n + \beta \cdot 3^n.
\]

\textbf{Քայլ 4. Կիրառել սկզբնական պայմանները։}

Տեղադրելով $n = 1$ և $n = 2$՝
\[
  \begin{cases}
    a_1 = 2\alpha + 3\beta = 2 \\[3pt]
    a_2 = 4\alpha + 9\beta = 10
  \end{cases}
\]

\textbf{Նշում.}\; Սա $2 \times 2$ գծային հավասարումների համակարգ է։

Առաջին հավասարումից՝
\[
  \alpha = \frac{2 - 3\beta}{2} = 1 - \frac{3\beta}{2}.
\]

Տեղադրելով երկրորդ հավասարման մեջ՝
\[
  4\!\left(1 - \frac{3\beta}{2}\right) + 9\beta = 10
  \quad\Longrightarrow\quad
  4 - 6\beta + 9\beta = 10
  \quad\Longrightarrow\quad
  3\beta = 6
  \quad\Longrightarrow\quad
  \beta = 2.
\]

Այդ դեպքում $\alpha = 1 - \dfrac{3 \cdot 2}{2} = 1 - 3 = -2$.

\textbf{Քայլ 5. Գրել մասնավոր լուծումը։}

\[
  a_n = -2 \cdot 2^n + 2 \cdot 3^n = \boxed{2\bigl(3^n - 2^n\bigr)}.
\]

\textbf{Քայլ 6. Ստուգել։}

\begin{center}
\renewcommand{\arraystretch}{1.3}
\begin{tabular}{cccc}
\toprule
$n$ & $2(3^n - 2^n)$ & Սպասվող & Համընկնում \\
\midrule
$1$ & $2(3 - 2) = 2$ & $a_1 = 2$ & $\checkmark$ \\
$2$ & $2(9 - 4) = 10$ & $a_2 = 10$ & $\checkmark$ \\
$3$ & $2(27 - 8) = 38$ & $5(10) - 6(2) = 38$ & $\checkmark$ \\
$4$ & $2(81 - 16) = 130$ & $5(38) - 6(10) = 130$ & $\checkmark$ \\
\bottomrule
\end{tabular}
\end{center}

\begin{answerbox}
\centering
$a_n = 2\bigl(3^n - 2^n\bigr)$
\end{answerbox}

% ══════════════════════════════════════════════════════════════════
\newpage
\section{Խնդիր 6 — Երկու Հաջորդական Զրո Պարունակող Բինար Տողեր}
% ══════════════════════════════════════════════════════════════════

\begin{theorybox}[Լրացման սկզբունք և անդրադարձ առնչություններ սահմանափակումներով]
\textbf{Հաշվարկ լրացման սկզբունքով։}\;
Երբեմն $P$ հատկությանը բավարարող օբյեկտների քանակը հաշվելիս ավելի հեշտ է լինում հաշվել այն օբյեկտները, որոնք \emph{չեն} բավարարում այդ պայմանին։
\[
  |P| = |\text{ընդհանուր}| - |\text{ոչ } P|\text{։}
\]

\textbf{Անդրադարձ առնչությունների կառուցում սահմանափակումներով տողերի համար։}\;
$n$ երկարությամբ բինար տողեր հաշվելու համար, որոնք բավարարում են (կամ չեն բավարարում) որևէ պայմանի՝
\begin{enumerate}[nosep]
  \item Սահմանենք հաջորդականություն. թող $b_n$-ը լինի $n$ երկարությամբ \emph{թույլատրելի} տողերի քանակը։
  \item \textbf{Դիտարկենք առաջին (կամ վերջին) նիշը։}\;
        Յուրաքանչյուր նիշի ընտրություն սահմանափակում է հաջորդող նիշերի տարբերակները՝ հանգեցնելով անդրադարձ առնչության։
  \item Լուծենք՝ օգտագործելով բազային դեպքերը։
\end{enumerate}

\medskip
\textbf{Օրինակ՝ տողեր առանց «$00$» ենթատողի։}\;
Դիտարկենք $n$ երկարությամբ տողեր, որտեղ չկան երկու հաջորդական $0$-ներ։
\begin{itemize}[nosep]
  \item Եթե տողը սկսվում է \texttt{1}-ով, ապա մնացած $n-1$ նիշերը կարող են կազմել ցանկացած թույլատրելի տող $\longrightarrow$ $b_{n-1}$ տարբերակ։
  \item Եթե տողը սկսվում է \texttt{0}-ով, ապա հաջորդ նիշը \emph{պարտադիր} պետք է լինի \texttt{1} (որպեսզի չլինի «$00$»), իսկ մնացած $n-2$ նիշերը կարող են կազմել ցանկացած թույլատրելի տող $\longrightarrow$ $b_{n-2}$ տարբերակ։
\end{itemize}
Սա տալիս է Ֆիբոնաչիի տիպի անդրադարձ առնչություն՝ $b_n = b_{n-1} + b_{n-2}$։

\begin{center}
\begin{tikzpicture}[>=Stealth, font=\small,
  level distance=2.0cm, sibling distance=5.0cm,
  edge from parent/.style={draw, thick, ->, >=Stealth},
  every node/.style={align=center}]
  \node[draw, rounded corners, fill=theoryBg, minimum width=2.0cm]
    {$n$ երկարության\\տող}
    child {
      node[draw, rounded corners, fill=setA!12, minimum width=2.0cm]
        {Սկսվում է \texttt{1}-ով\\մնացածը՝ $b_{n-1}$}
      edge from parent node[left=4pt, font=\footnotesize] {առաջին բիթ = \texttt{1}}
    }
    child {
      node[draw, rounded corners, fill=setB!12, minimum width=2.0cm]
        {Սկսվում է \texttt{01}-ով\\մնացածը՝ $b_{n-2}$}
      edge from parent node[right=4pt, font=\footnotesize] {առաջին բիթ = \texttt{0}}
    };
\end{tikzpicture}
\end{center}
\end{theorybox}

\begin{problembox}
10 երկարությամբ քանի՞ բինար տող կա, որ պարունակում է երկու հաջորդական զրո:
\end{problembox}

\subsection*{Լուծում}

Օգտագործենք \textbf{լրացման սկզբունքը}։

\textbf{Քայլ 1. Հաշվել 10 երկարությամբ բոլոր բինար տողերը։}

Յուրաքանչյուր դիրք ունի 2 հնարավորություն, ուստի ընդհանուր քանակը $2^{10} = 1024$ է։

\textbf{Քայլ 2. Հաշվել 10 երկարությամբ այն բինար տողերը, որոնք \emph{չեն պարունակում} երկու հաջորդական զրոներ։}

Թող $b_n$-ը լինի $n$ երկարությամբ այն բինար տողերի քանակը, որոնք չեն պարունակում երկու հաջորդական $0$-ներ։

\textbf{Պարզաբանում.}\; «Պարունակում է երկու հաջորդական զրո» նշանակում է, որ տողում կա \emph{առնվազն մեկ} «$00$» ենթատող։

\textbf{Առնչության դուրսբերումը։}\; $n = 3$-ի համար թույլատրելի (առանց «$00$») տողերն են՝
\texttt{010}, \texttt{011}, \texttt{101}, \texttt{110}, \texttt{111} — ճիշտ $5 = b_3$։
Ստուգում. \texttt{1}-ով սկսվող ցանկացած տող թողնում է 2 դիրք ($b_2 = 3$ թույլատրելի), իսկ \texttt{01}-ով սկսվող ցանկացած տող թողնում է 1 դիրք ($b_1 = 2$ թույլատրելի), տալով $3 + 2 = 5$։ \checkmark

\medskip
\textbf{Որոշումների ծառ $n = 3$-ի համար։}\;
Հետևյալ ծառը ցույց է տալիս 3-բիթանոց տողի բոլոր հնարավոր ընտրությունները՝ նշելով «$00$» պարունակող ճյուղերը որպես անթույլատրելի (\textcolor{setB}{\textsf{X}})։
Թույլատրելի տողերը (առանց «$00$») ցույց են տրված մնացած տերևներում։

\begin{center}
\begin{tikzpicture}[>=Stealth, font=\small,
  level 1/.style={sibling distance=6.4cm, level distance=1.6cm},
  level 2/.style={sibling distance=3.2cm, level distance=1.6cm},
  level 3/.style={sibling distance=1.6cm, level distance=1.6cm},
  edge from parent/.style={draw, thick, ->, >=Stealth},
  every node/.style={align=center},
  valid/.style={draw, rounded corners=3pt, fill=ansBg, minimum width=1.1cm,
                font=\small\ttfamily},
  invalid/.style={draw, rounded corners=3pt, fill=setB!15, minimum width=1.1cm,
                  font=\small\ttfamily, text=black!50},
  killed/.style={font=\bfseries\sffamily, text=setB}]
  \node[draw, rounded corners=4pt, fill=theoryBg] {սկիզբ}
    child { node[draw, rounded corners=3pt, fill=setA!10] {0\,\_\,\_}
      child { node[draw, rounded corners=3pt, fill=setB!15] {00\,\_}
        child { node[invalid] {000}
                edge from parent node[left, font=\scriptsize] {0} }
        child { node[invalid] {001}
                edge from parent node[right, font=\scriptsize] {1} }
        edge from parent node[left, font=\scriptsize] {0}
        node[killed, right=6pt, yshift=-2pt]
          {\textcolor{setB}{\textsf{X}\;\scriptsize անթույլատրելի}}
      }
      child { node[draw, rounded corners=3pt, fill=setA!10] {01\,\_}
        child { node[valid] {010}
                edge from parent node[left, font=\scriptsize] {0} }
        child { node[valid] {011}
                edge from parent node[right, font=\scriptsize] {1} }
        edge from parent node[right, font=\scriptsize] {1}
      }
      edge from parent node[left=3pt, font=\scriptsize] {0}
    }
    child { node[draw, rounded corners=3pt, fill=setA!10] {1\,\_\,\_}
      child { node[draw, rounded corners=3pt, fill=setA!10] {10\,\_}
        child { node[invalid] {100}
                edge from parent node[left, font=\scriptsize] {0}
                node[killed, right=3pt] {\textcolor{setB}{\textsf{X}}} }
        child { node[valid] {101}
                edge from parent node[right, font=\scriptsize] {1} }
        edge from parent node[left, font=\scriptsize] {0}
      }
      child { node[draw, rounded corners=3pt, fill=setA!10] {11\,\_}
        child { node[valid] {110}
                edge from parent node[left, font=\scriptsize] {0} }
        child { node[valid] {111}
                edge from parent node[right, font=\scriptsize] {1} }
        edge from parent node[right, font=\scriptsize] {1}
      }
      edge from parent node[right=3pt, font=\scriptsize] {1}
    };
\end{tikzpicture}
\end{center}

\noindent
\textbf{Մնացած տերևները} (կանաչ)՝ \texttt{010}, \texttt{011}, \texttt{101},
\texttt{110}, \texttt{111} — ճիշտ $\boxed{b_3 = 5}$, ինչը հաստատում է
$b_3 = b_2 + b_1 = 3 + 2 = 5$ առնչությունը։

\medskip
\textbf{Առնչության դուրսբերումը։}\;
Դիտարկենք $n$ երկարությամբ թույլատրելի տող.
\begin{itemize}[nosep]
  \item Եթե տողը սկսվում է \texttt{1}-ով, ապա մնացած $n-1$ նիշերը կարող են կազմել $n-1$ երկարությամբ ցանկացած թույլատրելի տող՝ $b_{n-1}$ տարբերակ։
  \item Եթե տողը սկսվում է \texttt{0}-ով, ապա հաջորդ նիշը պետք է լինի \texttt{1} (այլապես կունենանք «$00$»), և մնացած $n-2$ նիշերը կարող են կազմել $n-2$ երկարությամբ ցանկացած թույլատրելի տող՝ $b_{n-2}$ տարբերակ։
\end{itemize}

Հետևաբար՝
\[
  \boxed{b_n = b_{n-1} + b_{n-2}}, \qquad n \ge 3.
\]

\textbf{Բազային դեպքեր.}
\begin{itemize}[nosep]
  \item $b_1 = 2$: տողերն են \texttt{0} և \texttt{1}։
  \item $b_2 = 3$: թույլատրելի տողերն են \texttt{01}, \texttt{10}, \texttt{11}
        (բացառելով \texttt{00}-ն)։
\end{itemize}

\textbf{Քայլ 3. Հաշվել $b_n$-ը $n = 3, 4, \ldots, 10$ արժեքների համար։}

\begin{center}
\renewcommand{\arraystretch}{1.3}
\begin{tabular}{c c l}
\toprule
$n$ & $b_n$ & \textbf{Ստուգում}\; ($b_{n-1} + b_{n-2}$) \\
\midrule
$1$  & $2$   & (բազային դեպք՝ \texttt{0}, \texttt{1}) \\
$2$  & $3$   & (բազային դեպք՝ \texttt{01}, \texttt{10}, \texttt{11}) \\
\midrule
$3$  & $5$   & $b_2 + b_1 = 3 + 2 = 5$\;\checkmark \\
$4$  & $8$   & $b_3 + b_2 = 5 + 3 = 8$\;\checkmark \\
$5$  & $13$  & $b_4 + b_3 = 8 + 5 = 13$\;\checkmark \\
$6$  & $21$  & $b_5 + b_4 = 13 + 8 = 21$\;\checkmark \\
$7$  & $34$  & $b_6 + b_5 = 21 + 13 = 34$\;\checkmark \\
$8$  & $55$  & $b_7 + b_6 = 34 + 21 = 55$\;\checkmark \\
$9$  & $89$  & $b_8 + b_7 = 55 + 34 = 89$\;\checkmark \\
$10$ & $144$ & $b_9 + b_8 = 89 + 55 = 144$\;\checkmark \\
\bottomrule
\end{tabular}
\end{center}

\textbf{Դիտողություն.}\; $b_n = 2, 3, 5, 8, 13, 21, 34, 55, 89, 144$ արժեքները
Ֆիբոնաչիի թվերն են՝ $F_3, F_4, \ldots, F_{12}$։
Ընդհանուր դեպքում՝
\[
  \boxed{b_n = F_{n+2}}, \qquad \text{որտեղ } F_1 = 1,\; F_2 = 1,\; F_3 = 2,\; \ldots
\]

\textbf{Քայլ 4. Կիրառել լրացման սկզբունքը։}

\[
  \text{«$00$» պարունակող} = \text{Ընդհանուր} - \text{«$00$» չպարունակող}
  = 1024 - 144 = \boxed{880}.
\]

\textbf{Ստուգում փոքր արժեքների համար.}\;
$n = 2$-ի համար. ընդհանուր $= 4$, «$00$» չպարունակող $= 3$ (\texttt{01, 10, 11}),
ուստի «$00$» պարունակող $= 1$ (միայն \texttt{00})։ \checkmark

$n = 3$-ի համար. ընդհանուր $= 8$, $b_3 = 5$ (\texttt{010, 011, 101, 110, 111}),
ուստի «$00$» պարունակող $= 3$ (\texttt{001, 100, 000})։ \checkmark

\begin{answerbox}
\centering
10 երկարությամբ բինար տողերի քանակը, որոնք պարունակում են երկու հաջորդական զրո, \(\boxed{880}\) է։
\end{answerbox}

% ══════════════════════════════════════════════════════════════════
\newpage
\section{Խնդիր Հովսեփոսի մասին (Josephus Problem)}
% ══════════════════════════════════════════════════════════════════

\begin{problembox}
Գտեք $J(n)$ անդրադարձ առնչությունը, որը նկարագրում է $n$ մարդկանց շրջանից վերջին մնացողի համարը, եթե ամեն երկրորդ մարդը դուրս է գալիս։
\end{problembox}

\begin{tcolorbox}[colback=white, colframe=gray!50, arc=2mm, width=0.9\textwidth, halign=center]
  \textbf{Տեսանյութ.} Յոզեֆուսի խնդիրը (The Josephus Problem) \\
  \href{https://youtu.be/perxuYhrbzM}{\textbf{[VIDEO]} Դիտել YouTube-ում (ID: perxuYhrbzM)}
\end{tcolorbox}

\end{document}